\section{Implementation, evidence, and remaining limits}
\label{sec:implementation}

\subsection{The exact trusted boundary}

The selected checker is
\file{goal_checker/code/full_expression_checker.py}.
The readable block-form variant is
\file{goal_checker/code/full_checker.py}.
Each is a standalone Python module importing only \file{json} and
\file{sys}. No proof translator, external theorem database, generated
table, local helper, or previously checked proof is imported.

The input is an object with exactly four keys:
\file{variables}, \file{context}, \file{goal}, and \file{proof}.
The numerical scope is an exact nonnegative integer; Boolean indices
are rejected. Terms and formulas are tagged finite lists, with de Bruijn
indices. An iterator adds one bound variable in its update only;
$Q$ adds one and $D$ adds two in their formula bodies.

One recursive traversal checks constructor arities and child sorts,
performs variable actions, and normalizes exposed computations.
Validation calls it with the scoped identity action and requires its
output to equal the supplied syntax. Thus all supplied used terms
are canonical. The traversal checks \emph{all original children}
before erasing a computation, preventing an invalid index or wrong
sort from being hidden in a discarded body.

The checker validates a current goal before closing it as a retained
hypothesis or a reflexive equation. Those automatic leaves ignore
their proof payload; \file{null} is the conventional spelling.
This is an intentional input convention shared by the matched controls,
not a claim that unused JSON data have passed a separate proof-tree
grammar. All retained context formulas are validated, including
unused context entries. Non-leaf packets implement Tr, QE, DE,
or WI, checking every used major and parameter before substitution.

Success means exit status zero and exactly \file{ACCEPT} followed
by a newline. Every exception or exhausted resource is rejection,
not acceptance. Mathematical soundness is induction on the finite
accepting call tree; completeness is induction on a derivation,
supplying its intermediate major formulas in the packets.
The numerical interpretation is the ideal unbounded finite-tree
specification. Physical recursion and memory limits can reject
otherwise valid very large certificates.

\subsection{Reproducible measurements}

Each complete source is parsed with CPython 3.12.14 and measured after
\[
 \texttt{ast.unparse(ast.parse(source)) + '\textbackslash n'}.
\]
No AST nodes, docstrings, validation, normalization, or command-line
acceptance code are removed. Significant tokens exclude only comments
and structural whitespace tokens. AST counts include every node
returned by \file{ast.walk}. Deterministic gzip uses level $9$ and
timestamp zero; the recorded zlib runtime is 1.3.2.

The reported comparison is reproducible with
\begin{lstlisting}
python3 goal_checker/reproduce.py
python3 forward/verify.py
\end{lstlisting}
The first script regenerates the four whole-module measurements,
normalized sources, gzip streams, and its checker receipt.
The second regenerates the forward translation and normalization tests.
Neither script is part of a proof kernel. Other Python or zlib versions
may change normalized formatting or compressed lengths.

The large difference between the earlier $3095$-byte forward source and
the $2529$-byte selected checker is $566$ bytes, about $18.29\%$.
The packet format and validation organization changed, so this is a
whole-implementation comparison, not an isolated axiom experiment.
The matched table in Section~\ref{sec:result} is the appropriate test
of the incremental full-normalization change.

\subsection{Executed checks}

The goal-directed family passed $680$ command-line executions,
half in ordinary Python and half with \file{-O}: $440$ expected
acceptances and $240$ expected rejections, with no mismatches.
The $166$ distinct fixtures include every primitive rule, open
contexts, nested binders, scope escapes, Boolean indices, wrong
tags and arities, invalid data inside erased redexes, and bad
inference endpoints. Additional independent-forward comparisons
passed $500$ normalization checks, $500$ lifting checks, $250$
substitution checks, and $500$ canonical-validation checks.

The portable forward package translates all $64$ packaged
zero-normal source certificates. It passed $132$ command-line
checks, $1000$ normalization-idempotence checks, $1000$
substitution-coherence checks, $1000$ independent numerical
evaluation comparisons, and $2000$ defining-equation checks.
These are finite implementation tests. They supplement the
universal written proofs and are not a proof-assistant certificate.

One regression is particularly informative. The formula
\[
 D_{x,y}(y=1)
\]
is intuitionistically true under its reading: its conditional
antecedent is $0=1$, from which every $y=1$ follows.
A raw atomic \emph{case} introduction can prove it using a
nonliteral zero witness. That witness becomes literal $0$
after full normalization, where the same unmodified packet
would instead ask for progress from $z=1$ to $Sz=1$.
The naive translated certificate is rejected.
The compound-envelope translation produces an accepted
certificate for exactly the desired closed formula.
This test exercises a real introduction-mode mismatch.

\subsection{The external normalization cost}

Normalization is not a free oracle. Its executable recursion
is part of the kernel and its mathematical termination proof
is part of this report. It can nevertheless perform very
large computations and produce very large terms. The uniform
normalizer is not primitive recursive, as proved in
Section~\ref{sec:fullnf-metatheory}. No primitive-recursive bound
on this algorithm, and no internal iPRA proof of its uniform
totality, is asserted.

If one additionally requires a manifestly primitive-recursive
certificate relation, use explicit bounded computation traces
or the raw native calculus. A trace checker would be a different
measured implementation; its code is not silently omitted from
the present counts. The current user specification concerns
the size of a straightforward proof checker, so a total recursive
normalizing implementation is a legitimate candidate, with
this cost made explicit.

\subsection{What remains unresolved}

The result does not prove a lower bound on the size of any
iPRA checker, nor a universally simplest logical vocabulary.
Deleting an explicit computation axiom while implementing its
normalization is a change of proof presentation, not an
independence result showing that arithmetic computation has
ceased to be necessary. Reflexivity still seeds equality
reasoning, and the displayed mathematical inventory has six
patterns even though the program checks two of them as
automatic leaves.

For the requested complexity objective the evidence is mixed:
the selected matched source is shorter by three bytes, while
tokens, AST nodes, and gzip length increase. The stronger
result is the complete normalization and interpretation theorem,
including the compound-envelope repair. A substantial reduction
below all current checker profiles remains open.
